% GENERATED FILE. Edit canonical lessons, then run npm run generate:books.
% canonical-source-hash: fnv1a64:65618fee80f0fa19
% canonical-lessons: GE-C368-stattdessen, GE-C368-anstatt, GE-C368-ohne-dass, GE-C368-verzichten, GE-C368-ersetzen

\chapter{Stattdessen}
\label{ch:ge-stattdessen}
\hlchaptermodality{368}

\begin{chapteropening}
\textbf{By the end of this chapter:} \emph{I can say what happens instead, do without something and replace one thing with another.}

Complete the last lesson of chapter 368: I can say what happens instead, do without something and replace one thing with another.
\end{chapteropening}

\section[stattdessen]{stattdessen — instead}

\label{lesson:GE-C368-stattdessen}



\begin{quote}
Before the new one: say the German for after all, then the German for perhaps even.
\end{quote}

\subsection*{You'll want to know: stattdessen}
\textbf{stattdessen} — \textquotedblleft{}instead\textquotedblright{}.

\begin{quote}
Der Film war ausverkauft. Stattdessen sind wir essen gegangen.
\end{quote}

\begin{quote}
\emph{The film was sold out. Instead, we went out to eat.}
\end{quote}

The verb comes right after it: \emph{stattdessen sind wir …}.

\begin{cousinweb}
\emph{stattdessen} joins \emph{statt} \textquotedblleft{}instead of\textquotedblright{}, from \emph{Statt} \textquotedblleft{}place\textquotedblright{}, Old High German \emph{stat}, to \emph{dessen} \textquotedblleft{}of that\textquotedblright{}: in the place of that. EtymDB has no entry for the compound.
\end{cousinweb}

\begin{grammarlens}[title={What you've built}]
One more word for telling, explaining and describing at B1.
\end{grammarlens}

\subsection*{Your turn}
\begin{itemize}
\raggedright
  \item \emph{Say it:} \emph{stattdessen}
  \item \emph{Say it:} \emph{stattdessen}, once more
  \item \emph{Say it:} a sentence of your own with \emph{stattdessen}
  \item \emph{Recall:} say the German for after all, then the German for perhaps even, then say \emph{stattdessen} again
\end{itemize}

\begin{tcolorbox}[breakable,colback=teal!4,colframe=teal!35!black,title={Before you move on}]
What does stattdessen mean? (\textquotedblleft{}Instead\textquotedblright{}.) Say it once more.
\end{tcolorbox}
\section[anstatt]{anstatt — instead of}

\label{lesson:GE-C368-anstatt}



\begin{quote}
Before the new one: say the German for perhaps even, then the German for instead.
\end{quote}

\subsection*{You'll want to know: anstatt}
\textbf{anstatt} — \textquotedblleft{}instead of\textquotedblright{}.

\begin{quote}
Anstatt zu lernen, hat er den ganzen Abend gespielt.
\end{quote}

\begin{quote}
\emph{Instead of studying, he played all evening.}
\end{quote}

With \emph{zu} + infinitive, as here, or with the genitive: \emph{anstatt eines Briefes}. Also \emph{anstatt dass …}, with the verb at the end.

\begin{cousinweb}
\emph{anstatt} joins \emph{an} \textquotedblleft{}at, in\textquotedblright{} to \emph{Statt} \textquotedblleft{}place\textquotedblright{}, Old High German \emph{stat}: in the place of something else. EtymDB has no entry for it.
\end{cousinweb}

\begin{grammarlens}[title={What you've built}]
One more word for telling, explaining and describing at B1.
\end{grammarlens}

\subsection*{Your turn}
\begin{itemize}
\raggedright
  \item \emph{Say it:} \emph{anstatt}
  \item \emph{Say it:} \emph{anstatt}, once more
  \item \emph{Say it:} a sentence of your own with \emph{anstatt}
  \item \emph{Recall:} say the German for perhaps even, then the German for instead, then say \emph{anstatt} again
\end{itemize}

\begin{tcolorbox}[breakable,colback=teal!4,colframe=teal!35!black,title={Before you move on}]
What does anstatt mean? (\textquotedblleft{}Instead of\textquotedblright{}.) Say it once more.
\end{tcolorbox}
\section[ohne dass]{ohne dass — without (someone) doing}

\label{lesson:GE-C368-ohne-dass}



\begin{quote}
Before the new one: say the German for instead, then the German for instead of.
\end{quote}

\subsection*{You'll want to know: ohne dass}
\textbf{ohne dass} — \textquotedblleft{}without (someone) doing\textquotedblright{}.

\begin{quote}
Er ist gegangen, ohne dass jemand es gemerkt hat.
\end{quote}

\begin{quote}
\emph{He left without anyone noticing.}
\end{quote}

\emph{Ohne dass} sends the verb to the end. When the subject stays the same, German says \emph{ohne … zu}: \emph{Er ging, ohne etwas zu sagen.}

\begin{cousinweb}
\emph{ohne dass} joins \emph{ohne} \textquotedblleft{}without\textquotedblright{}, from Proto-Germanic \emph{ēnu} (EtymDB), a cousin of Greek \emph{áneu} \textquotedblleft{}without\textquotedblright{}, to \emph{dass} \textquotedblleft{}that\textquotedblright{}.
\end{cousinweb}

\begin{grammarlens}[title={What you've built}]
One more word for telling, explaining and describing at B1.
\end{grammarlens}

\subsection*{Your turn}
\begin{itemize}
\raggedright
  \item \emph{Say it:} \emph{ohne dass}
  \item \emph{Say it:} \emph{ohne dass}, once more
  \item \emph{Say it:} a sentence of your own with \emph{ohne dass}
  \item \emph{Recall:} say the German for instead, then the German for instead of, then say \emph{ohne dass} again
\end{itemize}

\begin{tcolorbox}[breakable,colback=teal!4,colframe=teal!35!black,title={Before you move on}]
What does ohne dass mean? (\textquotedblleft{}Without (someone) doing\textquotedblright{}.) Say it once more.
\end{tcolorbox}
\section[verzichten]{verzichten — to do without, to give up}

\label{lesson:GE-C368-verzichten}



\begin{quote}
Before the new one: say the German for instead of, then the German for without.
\end{quote}

\subsection*{You'll want to know: verzichten}
\textbf{verzichten} — \textquotedblleft{}to do without, to give up\textquotedblright{}.

\begin{quote}
Im Urlaub verzichte ich auf mein Handy.
\end{quote}

\begin{quote}
\emph{On holiday I do without my mobile.}
\end{quote}

\emph{Auf etwas verzichten} — to do without something. Inseparable: \emph{sie hat verzichtet}.

\begin{cousinweb}
\emph{verzichten} goes back to Middle High German \emph{verzihen} \textquotedblleft{}to renounce, to deny\textquotedblright{}. EtymDB does not give a usable history, so this follows the usual account.
\end{cousinweb}

\begin{grammarlens}[title={What you've built}]
One more word for telling, explaining and describing at B1.
\end{grammarlens}

\subsection*{Your turn}
\begin{itemize}
\raggedright
  \item \emph{Say it:} \emph{verzichten}
  \item \emph{Say it:} \emph{verzichten}, once more
  \item \emph{Say it:} a sentence of your own with \emph{verzichten}
  \item \emph{Recall:} say the German for instead of, then the German for without, then say \emph{verzichten} again
\end{itemize}

\begin{tcolorbox}[breakable,colback=teal!4,colframe=teal!35!black,title={Before you move on}]
What does verzichten mean? (\textquotedblleft{}To do without, to give up\textquotedblright{}.) Say it once more.
\end{tcolorbox}
\section[ersetzen]{ersetzen — to replace; to make up for}

\label{lesson:GE-C368-ersetzen}



\begin{quote}
Before the new one: say the German for without, then the German for to do without.
\end{quote}

\subsection*{You'll want to know: ersetzen}
\textbf{ersetzen} — \textquotedblleft{}to replace; to make up for\textquotedblright{}.

\begin{quote}
Die alte Lehrerin wird von einem jungen Kollegen ersetzt.
\end{quote}

\begin{quote}
\emph{The old teacher is being replaced by a young colleague.}
\end{quote}

Inseparable: \emph{er ersetzt, er hat ersetzt}. \emph{Den Schaden ersetzen} — to pay for the damage.

\begin{cousinweb}
\emph{ersetzen} is inherited from Middle High German \emph{ersetzen} (EtymDB), built with \emph{er-} on \emph{setzen} \textquotedblleft{}to set\textquotedblright{}, Proto-Germanic \emph{satjaną}: you set something new in the old place.
\end{cousinweb}

\begin{grammarlens}[title={What you've built}]
One more word for telling, explaining and describing at B1.
\end{grammarlens}

\subsection*{Your turn}
\begin{itemize}
\raggedright
  \item \emph{Say it:} \emph{ersetzen}
  \item \emph{Say it:} \emph{ersetzen}, once more
  \item \emph{Say it:} a sentence of your own with \emph{ersetzen}
  \item \emph{Recall:} say the German for without, then the German for to do without, then say \emph{ersetzen} again
\end{itemize}

\begin{tcolorbox}[breakable,colback=teal!4,colframe=teal!35!black,title={Before you move on}]
What does ersetzen mean? (\textquotedblleft{}To replace; to make up for\textquotedblright{}.) Say it once more.
\end{tcolorbox}
